\documentclass[a4paper,12pt,table]{article}

% The geometry package is used to customize page geometry
\usepackage[a4paper]{geometry}
% The tikz package supports custom graphics, such as the front page
\usepackage{tikz}	
% The babel package implements hyphenation and other typographical rules
\usepackage[english]{babel}
% The xcolor package provides enhanced support for color customization
\usepackage{xcolor}
% The graphicx package provides enhanced support for graphics
\usepackage{graphicx}
% The anyfontsize package is used for implementing unusual font sizes
\usepackage{anyfontsize}
% The fancyhdr package provides support for more complex headers and footers
\usepackage{fancyhdr}
% The caption package supports custom captioning of tables and figures
\usepackage{caption}
% The longtable package supports multi-page tables (longtabu)
\usepackage{longtable}
% The float package improves float placement (images, tables) and [H] specifier
\usepackage{float}
% The pifont package provides use of special characters (checkmarks, dings)
\usepackage{pifont}
% The pgfplots package provides support for tikz plots
\usepackage{pgfplots}
% The amsmath package provides support for page numbering
\usepackage{amsmath}
% The hyperref package provides support for reference links in the PDF
\usepackage{hyperref}
% The tabu package provides support for automatic text wrapping in table cells
\usepackage{tabu}
% The appendix package is for displaying the appendices correctly
\usepackage[titletoc]{appendix}
% The rotating packag is for being able to rotate images
\usepackage{rotating}

%Define normal page geometry
\newgeometry{layout=a4paper,hmargin=2cm,top=4.5cm,bottom=3.5cm,layoutoffset=0cm,footskip=17mm,headheight=2cm,headsep=15mm}

% The dateplot library provides support for dates on pgfplot axes
\pgfplotsset{compat=newest}
\usepgfplotslibrary{dateplot}
\usepgfplotslibrary{fillbetween}
\usepgfplotslibrary{groupplots}
\usepgfplotslibrary{colormaps}
% To save memory and compile time, externalize each plot
\usepgfplotslibrary{external}
\usetikzlibrary{external}
\tikzexternalize[optimize=false,prefix=auxiliary/,up to date check=diff]

% This is to use a single caption for multiple images in \center environment
 \graphicspath{ {figures/} }

% Set paragraph indentation and spacing
\setlength{\parindent}{0cm}		% Paragraph indentation
\setlength{\parskip}{3mm}			% Paragraph spacing

% Set table spacing parameters
\setlength\LTleft{0pt}			% Space between left margin and table
\setlength\LTright{0pt}			% Space between right margin and table
\tabulinesep=1.2mm					% Buffer for table cells

% Set font to Sans Serif for whole document (comment out for roman font)
\renewcommand{\familydefault}{\sfdefault}

% Set color palette for template
\definecolor{main}{HTML}{C64D00}
\definecolor{gray}{HTML}{404040}

% Define all new commands (standard, don't edit these)
\newcommand*{\Appendixautorefname}{Appendix}
\addto\extrasenglish{
	\def\tableautorefname{Table}
	\def\figureautorefname{Figure}
	\def\sectionautorefname{Section}
	\def\subsectionautorefname{Subsection}
}
\newcommand*{\tabref}[1]{\autoref{tab:#1}}
\newcommand*{\figref}[1]{\autoref{fig:#1}}
\newcommand*{\secref}[1]{\autoref{sec:#1}}
\newcommand*{\subsecref}[1]{\autoref{subsec:#1}}
\newcommand*{\appref}[1]{\autoref{app:#1}}
\newcommand*{\secauth}[1]{\vspace{-5mm}\hspace{10mm}\textit{~-~#1}\par}

% Set reference and number format
\setcounter{secnumdepth}{3}			% Number down to subsection
\setcounter{tocdepth}{3}				% Include down to subsection in TOC
\numberwithin{figure}{section}	% Number figures within subsections
\numberwithin{table}{section}		% Number tables within subsections

% Set up fancy header and footer
\pagestyle{fancy}		% Set all pages to fancy (except title)

% Header settings
\fancyhead{} %Clear all header fields
\fancyheadoffset{8mm}
% Logo in top left corner
\lhead{\includegraphics[height=1.8cm]{images/logo.png}}
% Report info in top right corner
\rhead{\color{gray} \textbf{\location\linebreak \subject\linebreak \doctitle}}

% Footer settings
\fancyfoot{} %Clear all footer fields
\fancyfootoffset{8mm}
% Three dots with page number on the bottom right corner of the page
\rfoot{
	\begin{tikzexternaldisable}
		\begin{tikzpicture}[remember picture, overlay]
			\fill[color=main](-47mm,0mm) circle (7mm);
			\fill[color=main!50](-27mm,0mm) circle (7mm);
			\fill[color=main!20](-7mm,0mm) circle (7mm);
			\node[color=gray] at (-7mm,0mm) {\textbf{\thepage}};
		\end{tikzpicture}
	\end{tikzexternaldisable}
}
% Classification info on the bottom left corner of the page
\lfoot{\vspace{-2mm}\color{gray} \textbf{CLASSIFICATION:\linebreak\classification}}

\fancypagestyle{plain}{
	% Footer settings
	\fancyfoot{} %Clear all footer fields
	\fancyfootoffset{-12mm}
	% Three dots with page number on the bottom right corner of the page
	\rfoot{
		\begin{tikzexternaldisable}
			\begin{tikzpicture}[remember picture, overlay]
				\fill[color=main](-47mm,0mm) circle (7mm);
				\fill[color=main!50](-27mm,0mm) circle (7mm);
				\fill[color=main!20](-7mm,0mm) circle (7mm);
			\end{tikzpicture}
		\end{tikzexternaldisable}
	}
	% Classification info on the bottom left corner of the page
	\lfoot{\vspace{-2mm}\color{gray} \textbf{CLASSIFICATION:\linebreak\classification}}
}